\phantomsection
\addcontentsline{toc}{subsection}{Restricción final}
\subsection*{Restricción final}

% Seccion con nombres de clase largos, que no se parten con guion.
\sloppy

\textit{El Facade debe poder configurarse desde un archivo YAML para definir secuencias de
actividades.}

\vspace{0.4cm}

\subsubsection*{Por qué esta restricción es la prueba de fuego del diseño}

Una secuencia que vive en un archivo solo es posible si la actividad es un \emph{dato}. Con la
fachada del inciso (b), donde cada actividad era un método con su cuerpo escrito, esta
restricción no tenía por dónde empezar: no hay forma de que un archivo añada un método a una
clase ya compilada.

El inciso (c) hizo el trabajo decisivo al convertir la actividad en un objeto. Aquí solo queda
construir ese objeto desde otra fuente, y por eso la clase que lo representa implementa la
misma interfaz:

{\small
\begin{verbatim}
public sealed class ActividadConfigurada : IActividad
{
    private readonly IReadOnlyList<PasoConfigurado> _pasos;
    private readonly IReadOnlyList<PasoConfigurado> _deshacer;

    public void Ejecutar()
    {
        foreach (PasoConfigurado paso in _pasos) { paso.Ejecutar(); }
    }
}
\end{verbatim}
}

Para la fachada no hay diferencia entre una actividad escrita en C\# y una definida en YAML.
Esa indiferencia es la restricción resuelta. También es la confirmación del argumento del
inciso (a) contra el método estático: una secuencia escrita en el cuerpo de un método estático
no se puede cargar de ninguna parte.

\subsubsection*{El archivo}

{\small
\begin{verbatim}
actividades:

  noche-de-cine:
    pasos:
      - luces.atenuar: 15
      - pantalla.bajar
      - proyector.encender
      - proyector.entrada: dvd
      - proyector.formato: 16:9
      - amplificador.encender
      - amplificador.entrada: dvd
      - amplificador.envolvente: si
      - amplificador.volumen: 6
      - dvd.encender
      - dvd.reproducir: El tercer hombre
    deshacer:
      - dvd.detener
      - dvd.apagar
      - amplificador.apagar
      - proyector.apagar
      - pantalla.subir
      - luces.encender
\end{verbatim}
}

Cada paso es \texttt{componente.accion} con un argumento opcional. La sección
\texttt{deshacer} no es un adorno: es lo que permite que una actividad leída de un archivo
cumpla el contrato completo de \texttt{IActividad}, incluido el \texttt{Deshacer} que el
inciso (c) introdujo.

\subsubsection*{Un lector escrito a mano, y por qué}

El ejercicio no tiene dependencias externas, así que el lector cubre a mano el subconjunto de
YAML que el archivo usa: mapas anidados por sangría de dos espacios, listas con guion,
escalares sin comillas y comentarios. No hay anclas, ni listas en línea, ni cadenas
multilínea, ni tipos.

Conviene decirlo claro: \textbf{en producción la elección correcta sería YamlDotNet}. Escribir
un analizador de un formato ajeno es una fuente conocida de sorpresas, y el valor de hacerlo
aquí es doble: no se añade una dependencia a un ejercicio de clase, y el subconjunto admitido
queda acotado y escrito en lugar de ser ``lo que la biblioteca acepte''.

\subsubsection*{La frontera entre el mundo tipado y el archivo}

Un archivo de configuración es, inevitablemente, cadenas. La pregunta de diseño es
\emph{dónde} se cruza esa frontera, y la respuesta elegida es un catálogo explícito:

{\small
\begin{verbatim}
public sealed class CatalogoAcciones
{
    private readonly Dictionary<string, Action<string?>> _acciones =
        new(StringComparer.OrdinalIgnoreCase);

    public CatalogoAcciones(Amplificador amplificador, ...)
    {
        _acciones["amplificador.encender"] = _ => amplificador.Encender();
        _acciones["amplificador.volumen"] =
            v => amplificador.PonerVolumen(Entero("amplificador.volumen", v));
        _acciones["luces.atenuar"] =
            v => luces.Atenuar(Entero("luces.atenuar", v));
        // ... veinte acciones en total
    }
}
\end{verbatim}
}

\textbf{Se construye a mano y sin reflexión, a propósito.} Resolver
\texttt{"amplificador.volumen"} con \texttt{Type.GetType} y \texttt{MethodInfo.Invoke} sería
más corto y bastante peor: cualquier nombre escrito en un archivo podría alcanzar cualquier
método de cualquier clase, incluidos los que nadie quiso exponer. El catálogo es una lista
blanca legible, y la guía metodológica del curso usa precisamente reflexión en el problema 1,
así que conviene nombrar el compromiso en lugar de copiarlo sin pensar.

\subsubsection*{Validar al cargar, no al ejecutar}

Es la decisión más importante de esta sección. El lector resuelve cada paso contra el catálogo
\emph{en el momento de cargar} y guarda la acción ya resuelta, no su nombre:

\begin{itemize}[leftmargin=1.2cm]
    \item Un paso inexistente hace fallar \textbf{la carga entera}, con el número de línea y
          el nombre que no existe.
    \item Se acumulan \emph{todos} los errores del archivo y se informan juntos, en vez de
          parar en el primero.
    \item Ninguno de los pasos válidos llega a ejecutarse. La alternativa —descubrir el error
          a mitad de la actividad— dejaría la sala a oscuras con el proyector encendido y la
          pantalla a medio bajar, que es exactamente el estado inconsistente que el usuario
          no sabe deshacer.
\end{itemize}

La parte 6 de la demostración lo comprueba con un archivo que pide
\texttt{amplificador.subir-mucho}.

\textbf{Y lo que queda sin comprobar.} El \emph{argumento} no se valida al cargar:
\texttt{luces.atenuar: muchisimo} pasa la carga y falla al ejecutarse, porque el catálogo
guarda un \texttt{Action<string?>} y la conversión ocurre dentro. Comprobar que la acción
existe es barato; comprobar que su argumento sirve exigiría que el catálogo declarara el tipo
esperado de cada acción. Es una limitación real, está dicha en el código y se arregla tipando
el catálogo si alguna vez importa.

\subsubsection*{Qué habilita}

\begin{itemize}[leftmargin=1.2cm]
    \item Un instalador puede definir ``noche de cine'' para un cliente concreto sin
          compilar nada.
    \item Las actividades dejan de ser código y pasan a ser configuración, con lo que eso
          implica: se versionan aparte, se revisan sin saber C\# y se cambian sin desplegar.
    \item Y el reverso, que conviene no olvidar: ahora hay un archivo más que puede estar mal,
          y la seguridad que antes daba el compilador la tiene que dar la validación de carga.
\end{itemize}

\subsubsection*{Evidencia de ejecución}

Salida literal de \texttt{DemoYaml}:

{\footnotesize
\begin{verbatim}
1. El archivo que define las secuencias
---------------------------------------
  actividades.yaml, copiado junto al ejecutable
  acciones invocables publicadas por el catalogo: 20

2. Cargar y registrar lo que diga el archivo
--------------------------------------------
  noche-de-cine       11 pasos
  tarde-de-consola    9 pasos

  la fachada no se recompilo ni se modifico para esto

3. Ejecutar una actividad que nadie escribio en C#
--------------------------------------------------
  cine.Ejecutar("noche-de-cine");
    luces         atenuadas al 15%
    pantalla      bajada
    proyector     encendido
    proyector     entrada -> dvd
    proyector     formato 16:9
    amplificador  encendido
    amplificador  entrada -> dvd
    amplificador  modo envolvente
    amplificador  volumen 6
    dvd           encendido
    dvd           reproduciendo "El tercer hombre"

4. Y deshacerla, porque el archivo tambien dice como
----------------------------------------------------
    dvd           detenido
    dvd           apagado
    amplificador  apagado
    proyector     apagado
    pantalla      subida
    luces         al 100%

5. La otra actividad del archivo
--------------------------------
  cine.Ejecutar("tarde-de-consola");
    luces         atenuadas al 65%
    pantalla      bajada
    proyector     encendido
    proyector     entrada -> consola
    amplificador  encendido
    amplificador  entrada -> consola
    amplificador  volumen 9
    consola       encendida
    consola       cargando "Celeste"
    consola       apagada
    amplificador  apagado
    proyector     apagado
    pantalla      subida
    luces         al 100%

6. Un archivo con errores falla al CARGAR, no a mitad de la sala
----------------------------------------------------------------
  el archivo pide "amplificador.subir-mucho", que no existe:
    El archivo de actividades tiene 1 problema(s):
      linea 5: no existe la accion "amplificador.subir-mucho"

  Ninguno de los tres pasos validos llego a ejecutarse: la sala
  no se quedo a oscuras con el proyector encendido.
  El valor "muchisimo" si pasa la carga y falla al ejecutarse:
  comprobar que existe la accion es barato, comprobar que el
  argumento sirve exigiria tipar el catalogo.
\end{verbatim}
}

\subsubsection*{Resumen del proceso}

El camino seguido en este ejercicio, que es el mismo que serviría para cualquier problema de
diseño:

\begin{enumerate}[leftmargin=1.2cm]
    \item Leer el enunciado buscando \emph{dónde está la dificultad}. Aquí no estaba en los
          componentes sino en el orden, que no vivía en ninguna clase.
    \item Escribir la solución ingenua y nombrar lo que rompe, sin caricaturizarla: con un
          cliente y una actividad habría sido la correcta.
    \item Elegir el patrón descartando candidatos con argumentos, incluidos los que se
          confunden con él, y responder expresamente a los dos que el enunciado nombra
          (inciso a).
    \item Implementar las tres actividades sobre los seis componentes de modo que cada
          actividad sea una coreografía distinta y no un parámetro (inciso b).
    \item Atacar la rigidez del resultado antes de que la pidan: Command para añadir
          actividades, Strategy para variar las existentes, midiendo el coste real de cada
          una (inciso c).
    \item Integrar los patrones que el enunciado pide diciendo lo que cuestan, no solo lo que
          dan: Observer salió casi gratis, Singleton no (inciso d).
    \item Cerrar con la restricción, que es la que comprueba si el diseño era de verdad
          extensible o solo lo parecía (esta sección).
\end{enumerate}

\fussy
